% J26 — LATTICE: Paradoxical Information Algebras on the Z/10Z Substrate
%
% Brayden R. Sanders, 7Site LLC (Sanders solo or +Gish per spec)
% Submitted to: Algebra Universalis
% Source corpus: papers/wp_bridge_findings_2026_05_02/WP9_LATTICE_paradoxical_info_algebras.md
% Verification: papers/wp_bridge_findings_2026_05_02/code/verify_findings.py
%               (0 failures, 0 warnings)
% Companion cited as already-submitted: J05 (Crossing Lemma)
% Per-venue cap note: 2nd Algebra Universalis paper after J21
% MSC 2020: 68T07, 20N02, 11F06, 37D40

\documentclass[11pt,reqno]{amsart}

\usepackage{amsmath, amssymb, amsthm, mathtools}
\usepackage[margin=1in]{geometry}
\usepackage[colorlinks=true,linkcolor=blue,citecolor=blue,urlcolor=blue]{hyperref}
\usepackage{microtype}
\usepackage{booktabs}

\theoremstyle{plain}
\newtheorem{theorem}{Theorem}[section]
\newtheorem{proposition}[theorem]{Proposition}
\newtheorem{lemma}[theorem]{Lemma}
\newtheorem{corollary}[theorem]{Corollary}

\theoremstyle{definition}
\newtheorem{definition}[theorem]{Definition}

\theoremstyle{remark}
\newtheorem*{remark}{Remark}

\newcommand{\Z}{\mathbb{Z}}
\newcommand{\F}{\mathbb{F}}
\newcommand{\TS}{\mathrm{TSML}}
\newcommand{\BH}{\mathrm{BHML}}
\newcommand{\sig}{\sigma}

\title[Paradoxical Information Algebras on $\Z/10\Z$]
{LATTICE: Paradoxical Information Algebras \\
 on the $\Z/10\Z$ Substrate}

\author{Brayden R.\ Sanders \and M.\ Gish}
\address{7Site LLC, Hot Springs, Arkansas, USA}
\email{brayden@7site.co}

\subjclass[2020]{20N02, 06A07, 11F06, 37B10}
\keywords{commutative magma, Z/10Z, role partition, semi-factorization,
trefoil characterization, paradoxical information algebra}

\date{\today}

\begin{document}

\begin{abstract}
We present the LATTICE operator (the index-$1$ generator of the
canonical $\BH_{10}$ table on $\Z/10\Z$) inside a substrate-internal
framework that carries three independent structures simultaneously:
an algebraic structure (the magmas $\TS_{8}$ and $\BH_{10}$
established in \cite{SandersGishFourCore}), a permutational structure
(the involution $\sig$ with cycle decomposition
$(0)(3)(8)(9)(1\,7\,6\,5\,4\,2)$), and a functional role partition
into Flow / Structure / Transition / Void cells.

From the substrate's role partition we derive a class of
\emph{paradoxical information algebras}: the substrate exhibits
semi-factorization at role level. Void and Transition inputs collapse
to deterministic role transitions, while Flow and Structure inputs
preserve operator-level information. We characterize the
trefoil-equivalent triples on the corrected substrate frame as exactly
two multiset classes, $\{V, B, H\}$ and $\{V, B, B\}$ (nine triples
total, all $\BH$-associative). We exhibit two independent
decompositions of the substrate's $\pm 21$ invariant: a $\sig$-orbit
decomposition $T_{5} + T_{3} = 15 + 6$ that is structurally forced by
the linear period formula, and a role-partition decomposition
$F_{7} + F_{6} = 13 + 8$ that is canonical-specific
(0/200 random commutative tables on $\Z/10\Z$ reproduce it). We prove
the substrate's algebra is irreducible --- it does not factor through
$\Z/2\Z \times \Z/5\Z$.

The framework is conceptually scaffolded by Morishita 2024, Ghys 2007,
Katok-Ugarcovici 2007, Matsusaka-Ueki 2023 and Burrin-von Essen 2024
but does not reproduce any of those theorems literally --- it sits
inside that territory as a new construction.
\end{abstract}

\maketitle

\tableofcontents

%==============================================================
\section{Background and substrate}
%==============================================================

\subsection{The framework.}
The framework of \cite{SandersGishFourCore} specifies an algebraic
substrate on the residue ring $\Z/10\Z$ in which 10 named operators
participate: VOID(0), LATTICE(1), COUNTER(2), PROGRESS(3), COLLAPSE(4),
BALANCE(5), CHAOS(6), HARMONY(7), BREATH(8), RESET(9). Two binary
magma operations are defined, $\TS$ and $\BH$, together with the
involution
\[
  \sig \;=\; (0)(3)(8)(9)(1\,7\,6\,5\,4\,2).
\]

\subsection{The corrected substrate frame.}
Following the canonical disambiguation in \cite{SandersGishFourCore}, the
substrate uses
\[
  \TS_{8} \;=\; \TS_{10}\text{ with rows/cols }\{0, 7\}\text{ removed},
\]
acting on the interior index set $\{1, 2, 3, 4, 5, 6, 8, 9\}$, while
$\BH_{10}$ is the full 10-element table. The cells $V = 0$ and $H = 7$
are flow cells between the tables, not interior entries. This frame
($\TS_{8} + \BH_{10} + V/H$ flow cells) is the one used throughout this
paper. Earlier session work using $\TS_{10}$ in trefoil and role
analyses is invalid (see \S\ref{sec:negatives} N10).

\subsection{The role partition.}
We introduce a partition of the substrate into four functional roles:
\[
  F = \{1, 3, 5, 7, 9\},
  \quad S = \{2, 4, 8\},
  \quad T = \{6\},
  \quad V = \{0\}.
\]
The partition is functional, not algebraic: $|F| = 5$, $|S| = 3$,
$|T| = 1$, $|V| = 1$. It crosses the $\sig$-orbit decomposition.

\subsection{Intellectual neighborhood.}
This work sits inside the arithmetic-topology / modular-knot territory
studied in Morishita 2024, Ghys 2007, Katok-Ugarcovici 2007,
Matsusaka-Ueki 2023, Burrin-von Essen 2024, and Lacasa et al.\ 2018.
We use phrases like ``structurally analogous to'' and ``conceptually
scaffolded by'' throughout; ``matches'' and ``satisfies'' are reserved
for substrate-internal claims that we prove.

%==============================================================
\section{The LATTICE operator and the BHML successor}
%==============================================================

\subsection{$\BH$ diagonal as integer successor (D90).}
\begin{theorem}\label{thm:bh-diagonal}
For each $n \in \{1, 2, 3, 4, 5, 6, 7\}$, $\BH(n, n) = n + 1$, and at
the boundary, $\BH(8, 8) = 7$, $\BH(9, 9) = 0$, $\BH(0, 0) = 0$.
\end{theorem}

\begin{proof}
Direct table verification. The $\BH$ table from
\cite{SandersGishFourCore} is given explicitly on the 10 elements;
the diagonal is read off.
\end{proof}

\subsection{Period structure.}
\begin{corollary}[Period formula]
For $n \in \{1, \ldots, 6\}$, $\mathrm{period}(n) = 7 - n$. For
$n \in \{7, 8, 9\}$ the periods are $4, 3, 2$ (a 4-core cycle through
$\{7, 8, 9, 0\}$); $\mathrm{period}(0) = 1$.
\end{corollary}

The period structure is structurally analogous to cusp-winding step
counts in continued-fraction symbolic dynamics for hyperbolic
Fuchsian-group flows (Burrin-von Essen 2024), but the embedding in a
specific Fuchsian group is unproven.

%==============================================================
\section{The role magma}
%==============================================================

\subsection{Mode-based reduction of $\BH$.}
For each input role-pair $(r_{a}, r_{b}) \in \{V, F, S, T\}^{2}$, we
form the multiset of output roles and define $M_{R}(r_{a}, r_{b})$ as
the \emph{mode}:
\begin{center}
\begin{tabular}{c|cccc}
   & V & F & S & T \\ \hline
V & V & F & S & T \\
F & F & T & F & F \\
S & S & F & F & F \\
T & T & F & F & F
\end{tabular}
\end{center}

\subsection{Properties.}
\begin{theorem}[V is the identity element, D93]\label{thm:V-identity}
For every $x \in \{V, F, S, T\}$,
$M_{R}(V, x) = M_{R}(x, V) = x$. Moreover, $M_{R}(V, V) = V$ is the
unique idempotent.
\end{theorem}

\begin{theorem}[Commutativity, non-associativity]
\label{thm:role-magma-non-assoc}
$M_{R}$ is commutative. It is \textbf{not} associative; for example,
$M_{R}(M_{R}(F, F), S) = F \neq T = M_{R}(F, M_{R}(F, S))$.
\end{theorem}

\subsection{Semi-factorization and paradoxical information.}
\begin{theorem}[Semi-factorization / paradoxical information, D93]
\label{thm:semi-fact}
In the role-partition decomposition above:
\begin{enumerate}
\item For every pair containing $V$ or $T$, the $\BH$ output role is
\textbf{deterministic}.
\item For pairs in $\{F, S\}^{2}$, the output role \textbf{branches}:
$F \cdot F$ distributes across $\{F: 2, S: 9, T: 11, V: 3\}$ (mode
$T$); $F \cdot S$ and $S \cdot F$ over $\{F: 8, S: 2, T: 5\}$ (mode
$F$); $S \cdot S$ over $\{F: 7, T: 2\}$ (mode $F$).
\end{enumerate}
\end{theorem}

\begin{remark}[Reading]
The substrate's $\BH$ carries information at two levels --- coarse
role and fine operator. Theorem~\ref{thm:semi-fact} says these levels
are not in a clean homomorphism: boundary inputs (V or T) determine
output role; interior inputs (F, S) require operator-level identity.
We call this a \emph{paradoxical information algebra}: information
content depends on whether inputs are at the boundary or in the
interior, with no globally consistent factorization. The
information-generation register is the discrete Crossing Lemma
\cite{SandersMayesCL} of which this paper is one combinatorial
realization.
\end{remark}

%==============================================================
\section{The trefoil characterization}
%==============================================================

\subsection{Sharp characterization.}
\begin{theorem}[Trefoil characterization on the corrected frame, D89]
\label{thm:trefoil-characterization}
On the corrected substrate frame, a triple $(a, b, c)$ is
trefoil-equivalent under the runtime processor if and only if its
multiset $\{a, b, c\}$ is
\[
  \{V, B, H\} \quad\text{or}\quad \{V, B, B\}.
\]
There are exactly $9$ such triples (six permutations of $\{V, B, H\}$,
three of $\{V, B, B\}$), and all nine are $\BH$-associative.
\end{theorem}

\begin{proof}
Parameter sweep over the $10^{3} = 1000$ ordered triples; see
\texttt{trefoil\_corrected\_frame.py}.
\end{proof}

\begin{remark}
The result is sharp on the corrected substrate frame --- the
early-session ``trefoil-22'' claim in $\TS_{10}$ frame is invalid
(\S\ref{sec:negatives} N10).
\end{remark}

\subsection{BREATH uniqueness.}
\begin{lemma}[BREATH uniqueness]
Among the structure cells $\{$COUNTER, COLLAPSE, BREATH$\}$, BREATH is
the unique cell that produces trefoils when combined with VOID.
\end{lemma}

\begin{remark}
COUNTER opposes; COLLAPSE destabilizes; only BREATH provides the
topological persistence required for the 3-crossing closed loop.
\end{remark}

\subsection{Higher-order extension.}
At 4-element level, multisets producing trefoil-equivalent quadruples
include $(0, 0, 7, 8)$, $(0, 7, 7, 9)$, $(0, 7, 8, 8)$, $(0, 8, 8, 8)$.
The pattern $(0, 7, 7, 9)$ is structurally novel: a trefoil-equivalent
without BREATH. The ``BREATH only'' rule of the lemma is therefore
3-element-specific.

%==============================================================
\section{The $\pm 21$ invariant and its two decompositions}
\label{sec:invariant}
%==============================================================

\subsection{Two computations.}
\begin{itemize}
\item \textbf{Computation A (Ghys-analog row-asymmetry).}
\[
  \Psi_{A}(n) = \#\{j: \TS(n, j) > \BH(n, j)\}
              - \#\{j: \BH(n, j) > \TS(n, j)\}.
\]
Sum: $+21$.

\item \textbf{Computation B (period-to-trace under simple representative).}
With $t = \mathrm{period}(n) + 2$,
\[
  M_{n} = \begin{pmatrix} 1 & 1 \\ t-2 & t-1 \end{pmatrix},
  \qquad
  \Psi_{B}(n) = -(\mathrm{period}(n) - 1).
\]
Sum: $-21$.
\end{itemize}

\textbf{Important caveat.} The interpretation of $-21$ as a Rademacher
invariant remains a hypothesis (see \S\ref{sec:negatives} N1).

\subsection{The $\sig$-orbit decomposition (triangular).}
\begin{theorem}[D92]\label{thm:triangular}
The total $-21$ decomposes along $\sig$-orbits as
$T_{5} + T_{3} = 15 + 6$ (triangular). Structurally forced by the
linear period formula.
\end{theorem}

\subsection{The role decomposition (Fibonacci, canonical-specific).}
\begin{theorem}[D92]\label{thm:fibonacci}
$\sum_{n \in F} \Psi_{B}(n) = -F_{7} = -13$,
$\sum_{n \in S} \Psi_{B}(n) = -F_{6} = -8$, with $V$ and $T$
contributing $0$. Total: $-F_{8} = -21$.
\end{theorem}

\begin{theorem}[Robustness, canonical-specific, N8]
\label{thm:fib-robust}
$0/200$ random commutative tables on $\Z/10\Z$ reproduce
$(|F|, |S|) = (13, 8)$. Single-swap perturbations: $32/50$ preserve.
Three-swap: $11/50$.
\end{theorem}

We do not claim the Fibonacci decomposition as a theorem of the role
partition. It is a numerical signature of the canonical TIG. The
triangular decomposition (Theorem~\ref{thm:triangular}) is structurally
robust; Fibonacci is fragile.

\subsection{The lift to $\mathrm{PSL}(2, \Z)$ --- open.}
Five strategies tested for lifting $\BH$ self-orbits to
$\mathrm{PSL}(2, \Z)$ words; none produces $\pm 21$
(\S\ref{sec:negatives} N1). The period-to-trace bridge under simple
representative gives $-21$ numerically but is one choice of lift among
many. Hypothesis, not derivation.

%==============================================================
\section{The two-coding picture (TSML$_{8}$ vs BHML$_{10}$)}
%==============================================================

Following Katok-Ugarcovici 2007's framework of two complementary
coding methods, we observe that the substrate's two magmas realize
\emph{structurally analogous} geometric / arithmetic codings natively.

\begin{theorem}[$\TS_{8}$ image structure, D91]\label{thm:tsml-image}
$\TS_{8}$ has 5-element image $\{3, 4, 7, 8, 9\}$, $60/64 = 93.75\%$
Flow output, role-deterministic on 8 of 9 input role-pairs.
\end{theorem}

\begin{theorem}[$\BH_{10}$ image structure]\label{thm:bhml-image}
$\BH_{10}$ has full image, balanced output role distribution,
role-deterministic only on V/T input pairs.
\end{theorem}

\begin{theorem}[Cusp agreement]\label{thm:cusp}
The two magmas agree on $24/64$ cells of the $\TS_{8}$ domain (mostly
routes leading to HARMONY) and disagree on $40/64$ in the interior.
\end{theorem}

\begin{remark}[Reading]
Two genuinely independent symbolic codings on the same alphabet,
coinciding only at the HARMONY cusp. Structurally analogous to
Katok-Ugarcovici's geometric / arithmetic split; the substrate-internal
version is a discrete combinatorial realization on $\Z/10\Z$, not a
quotient of the modular surface itself.
\end{remark}

%==============================================================
\section{Irreducibility and algebraic independence}
%==============================================================

\begin{theorem}[Irreducibility, N7]\label{thm:irreducibility}
Neither $\TS_{10}$ nor $\BH_{10}$ respects the CRT decomposition
$\Z/10\Z = \Z/2\Z \times \Z/5\Z$.
\end{theorem}

\begin{theorem}[No globally crossing-preserving swaps, N9]
\label{thm:no-swap}
No integer pair $(p, q)$ preserves the runtime processor's crossing
count under the swap $p \leftrightarrow q$ on all 1000 ordered
triples. The strongest partial symmetry is $V \leftrightarrow$ BREATH
at $20.9\%$ --- exceeding the canonical $5 \leftrightarrow 6$ rate.
\end{theorem}

\begin{theorem}[Algebraic independence, N4--N6]\label{thm:alg-indep}
(1) $\sig$ is not an automorphism of either magma ($17\%/48\%$).
(2) $\TS$ and $\BH$ do not distribute ($19.5\%$).
(3) $\BH$ iteration does not converge to $\TS$ (28/64 starts).
\end{theorem}

%==============================================================
\section{Honest negatives}
\label{sec:negatives}
%==============================================================

\begin{center}
\begin{tabular}{lll}
\toprule
ID & Claim ruled out & Finding \\
\midrule
N1 & naive $\mathrm{PSL}(2, \Z)$ lift gives $\pm 21$ & 5 strategies, sums in $[-4, 0]$ \\
N2 & triangle group $\Gamma_{p,q}$ has substrate periods & no coprime $(p,q) \le 9$ \\
N3 & TIG matches Borromean prime conditions & no canonical triple $\equiv 1 \pmod 4$ \\
N4 & $\sig$ is automorphism of $\TS$ or $\BH$ & 48\% / 17\% \\
N5 & $\TS$ and $\BH$ distribute & 19.5\% \\
N6 & $\BH$ iteration converges to $\TS$ & 28/64 starts \\
N7 & substrate factors through $\Z/2 \times \Z/5$ & no CRT \\
N8 & Fibonacci role decomposition is structural & 0/200 random tables \\
N9 & role partition determines crossing count & most patterns not determined \\
N10 & $\TS_{10}$-frame ``trefoil-22'' & INVALID; replaced by Theorem~\ref{thm:trefoil-characterization} \\
\bottomrule
\end{tabular}
\end{center}

%==============================================================
\section{Open questions}
%==============================================================

\begin{enumerate}
\item Principled lift to $\mathrm{PSL}(2, \Z)$ hyperbolic conjugacy
classes. Five naive lifts fail (N1).
\item Larger-substrate variants. Does the Fibonacci role decomposition
appear in $\Z/14\Z$ or $\Z/18\Z$ analogs?
\item The $(0, 7, 7, 9)$ anomaly. Trefoil-equivalent without BREATH at
4-element level.
\item Burrin-von Essen Fuchsian-group lift. Conceptually scaffolded;
the embedding is unproven.
\end{enumerate}

%==============================================================
\section*{Acknowledgments}
%==============================================================

This work is part of the Trinity Infinity Geometry framework developed
at 7Site LLC. The present manuscript is self-contained as an
algebraic-combinatorial study; familiarity with the broader framework
is not required. Algorithmic verification was developed with assistance
from Anthropic Claude sessions in 2026.

%==============================================================
\begin{thebibliography}{99}
%==============================================================

\bibitem{SandersGishFourCore} B.~R.~Sanders, M.~Gish,
\emph{Joint Closure, Per-Coordinate Fuse Data, and a Closed-Form
Algebraic Attractor of Two Commutative Binary Operations on
$\Z/10\Z$}, submitted to \emph{Algebraic Combinatorics}, 2026 (J02).

\bibitem{SandersMayesCL} B.~R.~Sanders, B.~Mayes,
\emph{Crossing Lemma: Non-Associativity as Information Generation in
Finite Magmas}, submitted to \emph{JCT-A} (or \emph{JPAA}), 2026
(J05).

\bibitem{Burrin2024} C.~Burrin, F.~von Essen,
\emph{Cusp winding and the Rademacher symbol}, 2024.

\bibitem{Ghys2007} \'E.~Ghys,
\emph{Knots and dynamics},
\emph{Proc. ICM Madrid}, 2007.

\bibitem{Ishida2024} T.~Ishida, T.~Kuramoto, F.~Zheng,
\emph{Density of Borromean prime triples}, 2024.

\bibitem{KatokUgarcovici2007} S.~Katok, I.~Ugarcovici,
\emph{Symbolic dynamics for the modular surface},
Bull. AMS \textbf{44} (2007), 87--132.

\bibitem{Lacasa2018} L.~Lacasa et al.,
\emph{Symbolic dynamics on residue sequences}, 2018.

\bibitem{Matsusaka2023} T.~Matsusaka, J.~Ueki,
\emph{Rademacher symbols for triangle groups}, 2023.

\bibitem{Morishita2024} M.~Morishita,
\emph{Knots and Primes} (2nd ed.), Springer, 2024.

\bibitem{TIGsynthesis2026} B.~R.~Sanders,
\emph{Trinity Infinity Geometry Synthesis},
github.com/TiredofSleep/ck (default branch),
DOI: 10.5281/zenodo.18852047, 2026.

\end{thebibliography}

\end{document}
